%============================================================
% CHAPITRE 3 — ANALYSE ET CONCEPTION
%============================================================
\chapter{Analyse et Conception}

\section{Introduction}

La phase d'analyse et de conception occupe une place centrale dans tout projet
de développement logiciel. Elle permet de traduire les exigences recueillies
au chapitre précédent en un modèle architectural cohérent, clair et techniquement
réalisable. Pour la modélisation de MediCab Pro, nous avons fait le choix du
langage UML (Unified Modeling Language), standard reconnu pour sa capacité à
représenter fidèlement les architectures distribuées complexes et à faciliter
la communication entre toutes les parties prenantes du projet.

\section{Architecture technique globale}

\subsection{Vue d'ensemble de l'architecture}

L'architecture de MediCab Pro repose sur un découpage en couches fonctionnelles
bien séparées. La couche de présentation, développée en Angular, ne communique
jamais directement avec les microservices applicatifs : toutes les requêtes
transitent par une API Gateway qui joue le rôle de point d'entrée centralisé,
chargé du routage, de la sécurité et de l'équilibrage de charge. Chaque
microservice Spring Boot expose ses propres fonctionnalités via une interface
REST documentée avec Swagger. La communication inter-services s'effectue de
manière synchrone via REST, renforcée par des mécanismes de tolérance aux
pannes (Circuit Breaker Resilience4j).

\begin{figure}[H]
  \centering
  \includegraphics[width=0.9\textwidth]{Architecture globale.png}
  \caption{Architecture microservices globale de MediCab Pro}
  \label{fig:arch-globale}
\end{figure}

\subsection{Microservices implémentés}

\begin{table}[H]
\centering
\caption{Inventaire des microservices de MediCab Pro}
\label{tab:microservices}
\renewcommand{\arraystretch}{1.4}
\begin{tabularx}{\textwidth}{|>{\bfseries\ttfamily}p{4cm}|p{1.5cm}|X|}
\hline
\textbf{Service} & \textbf{Port} & \textbf{Responsabilité} \\
\hline
auth-service & 8081 & Authentification JWT, OAuth2 et contrôle d'accès RBAC \\
\hline
appointment-service & 8082 & Gestion des agendas, rendez-vous et file d'attente \\
\hline
patient-service & 8083 & Dossiers patients, antécédents et documents médicaux \\
\hline
medical-record-service & 8084 & Consultations médicales et historique clinique \\
\hline
ordonnance-service & 8085 & Prescriptions électroniques et gestion médicaments \\
\hline
billing-service & 8086 & Facturation, suivi des paiements et rapports financiers \\
\hline
notification-service & 8087 & Envoi de notifications, alertes et rappels \\
\hline
ai-service & 8088 & Module IA MedAgent (FastAPI / Python) \\
\hline
api-gateway & 8080 & Point d'entrée unique, routage et sécurité transverse \\
\hline
discovery-service & 8761 & Registre de services Eureka Server \\
\hline
config-service & 8888 & Configuration externalisée et centralisée \\
\hline
\end{tabularx}
\end{table}

\section{Diagrammes de cas d'utilisation}

\subsection{Définition}

Le diagramme de cas d'utilisation matérialise les interactions entre les acteurs
externes au système et les fonctionnalités que celui-ci met à leur disposition.
Il constitue le premier niveau de modélisation UML et permet de capturer les
exigences fonctionnelles sous une forme visuelle, accessible à toutes les parties
prenantes, qu'elles soient techniques ou métier.

\subsection{DCU --- Espace Médecin}

\begin{figure}[H]
  \centering
  \includegraphics[width=0.9\textwidth]{UsecaseMedcine.png}
  \caption{Diagramme de cas d'utilisation --- Espace Médecin}
  \label{fig:dcu-medecin}
\end{figure}

\subsection{DCU --- Espace Secrétaire Médicale}

\begin{figure}[H]
  \centering
  \includegraphics[width=\textwidth]{UsecaseSecritaire.png}
  \caption{Diagramme de cas d'utilisation --- Espace Secrétaire Médicale}
  \label{fig:dcu-secretaire}
\end{figure}

\subsection{DCU --- Espace Super Administrateur}

\begin{figure}[H]
  \centering
  \includegraphics[width=\textwidth]{UsecaseSuperAdmin.png}
  \caption{Diagramme de cas d'utilisation --- Espace Super Administrateur}
  \label{fig:dcu-admin}
\end{figure}

\subsection{DCU --- Module IA (MedAgent)}

\begin{figure}[H]
  \centering
  \includegraphics[width=\textwidth]{UsecaseAiAgent.png}
  \caption{Diagramme de cas d'utilisation --- Module IA MedAgent}
  \label{fig:dcu-ia}
\end{figure}

\section{Diagrammes de séquence}

\subsection{Diagramme de séquence --- Authentification JWT}

\begin{figure}[H]
  \centering
  \includegraphics[width=\textwidth]{SequenceAuth.png}
  \caption{Diagramme de séquence --- Authentification JWT}
  \label{fig:ds-auth}
\end{figure}

\subsection{Diagramme de séquence --- Création d'un rendez-vous}

\begin{figure}[H]
  \centering
  \includegraphics[width=\textwidth]{SequanceRendezVous.png}
  \caption{Diagramme de séquence --- Création d'un rendez-vous}
  \label{fig:ds-rdv}
\end{figure}

\subsection{Diagramme de séquence --- Consultation avec MedAgent}

\begin{figure}[H]
  \centering
  \includegraphics[width=\textwidth]{SequanceConsultationAiAgent.png}
  \caption{Diagramme de séquence --- Consultation assistée par MedAgent}
  \label{fig:ds-medagent}
\end{figure}

\subsection{Diagramme de séquence --- Génération de facture}

\begin{figure}[H]
  \centering
  \includegraphics[width=\textwidth]{SequanceFacture.png}
  \caption{Diagramme de séquence --- Génération de facture}
  \label{fig:ds-facture}
\end{figure}

\section{Diagrammes de classes}

\subsection{Module Authentification (auth-service)}

\begin{figure}[H]
  \centering
  \includegraphics[width=\textwidth]{ClasseAuth.png}
  \caption{Diagramme de classe --- Module Authentification}
  \label{fig:dc-auth}
\end{figure}

\subsection{Module Patient (patient-service)}

\begin{figure}[H]
  \centering
  \includegraphics[width=\textwidth]{ClassePatient.png}
  \caption{Diagramme de classe --- Module Patient}
  \label{fig:dc-patient}
\end{figure}

\subsection{Module Rendez-vous (appointment-service)}

\begin{figure}[H]
  \centering
  \includegraphics[width=\textwidth]{ClasseRendesVous.png}
  \caption{Diagramme de classe --- Module Rendez-vous}
  \label{fig:dc-rdv}
\end{figure}

\subsection{Module Consultations et Ordonnances}

\begin{figure}[H]
  \centering
  \includegraphics[width=\textwidth]{ClasseOrdonanceConsultataion.png}
  \caption{Diagramme de classe --- Module Consultations et Ordonnances}
  \label{fig:dc-consultation}
\end{figure}

\subsection{Module Facturation (billing-service)}

\begin{figure}[H]
  \centering
  \includegraphics[width=\textwidth]{ClasseFacturation.png}
  \caption{Diagramme de classe --- Module Facturation}
  \label{fig:dc-billing}
\end{figure}

\section{Conclusion}

Ce chapitre a exposé l'ensemble de la conception architecturale de MediCab Pro :
une structure distribuée composée de onze microservices indépendants, une
modélisation UML complète des interactions entre acteurs et composants, et une
conception guidée par le principe de séparation des responsabilités. Le chapitre
suivant décrit la phase de réalisation concrète de l'application.
